%%%%%%%%%%%%%%%%%%%%%%%%
%
% $Autor: Gautam Ramesh$
% $Datum: 2025-10-14 21:52:19Z $
% $Pfad: /Users/gautamramesh/githubNew/ML26-03-Car-Licence-Plate-Identification-with-a-Jetson-Nano/Assembly - edit/Chapters/en/SafetyGuidelines.tex $
% $Version: 1 $
%
% !TeX encoding = utf8
% !TeX root = Rename
%
%%%%%%%%%%%%%%%%%%%%%%%%



\chapter{Safety Guidelines}

This section outlines essential safety procedures to prevent hardware damage, data loss, or user injury during operation, maintenance, and transportation of the Jetson Nano–based License Plate Detection System.  
Always follow these safety recommendations before connecting, testing, or relocating the unit.

% -------------------------------------------------------
\section{General Safety Instructions}

\begin{itemize}
	\item \textbf{Qualified Personnel Only:} Setup and maintenance should be performed by individuals familiar with embedded systems and electrical safety.
	\item \textbf{Disconnect Power:} Always unplug the \SI{5}{\volt}/\SI{4}{\ampere} power supply before connecting or removing peripherals.
	\item \textbf{Static Protection:} Use an anti-static wrist strap or work on an ESD-safe mat to prevent electrostatic discharge damage to sensitive components.
	\item \textbf{Avoid Liquid Contact:} Do not expose the system to moisture or conductive liquids. Keep the workspace dry and well ventilated.
	\item \textbf{Ventilation:} Operate the device on an open surface; ensure sufficient airflow around heatsinks and fans.
	\item \textbf{Stable Power Source:} Use only the certified NVIDIA Jetson Nano adapter or equivalent with correct polarity and amperage.
\end{itemize}

\vspace{0.4cm}

% -------------------------------------------------------
\section{Electrical and Operational Safety}

\begin{table}[H]
	\centering
	\caption{Electrical and Operational Safety Checklist}
	\begin{tabular}{|L{3.8cm}|L{8.2cm}|L{2.2cm}|}
		\hline
		\textbf{Item} & \textbf{Precaution / Action} & \textbf{Frequency} \\ \hline
		Power Adapter & Verify rated \SI{5}{\volt}, \SI{4}{\ampere}; check barrel connector for looseness or heat & Each setup \\ \hline
		USB Ports & Avoid hot-plugging high-current devices (e.g., external HDDs) without a powered hub & Always \\ \hline
		GPIO Header & Connect only 3.3V logic; avoid 5V direct inputs & Each wiring \\ \hline
		HDMI Cable & Insert gently; avoid bending near connector & Each setup \\ \hline
		Cooling System & Ensure fan rotation; do not block airflow during operation & Continuous \\ \hline
		Camera Module & Secure cable connection; protect from static and dust & Before each run \\ \hline
	\end{tabular}
\end{table}

These safety practices are designed to prevent hardware failure, port burnout, and accidental short circuits during extended operation.

\vspace{0.4cm}

% -------------------------------------------------------
\section{Environmental and Handling Guidelines}

\begin{itemize}
	\item \textbf{Temperature:} Operate between \SI{10}{\celsius} and \SI{35}{\celsius}. Avoid direct sunlight or high ambient heat sources.
	\item \textbf{Humidity:} Keep humidity below 70\%; condensation can cause corrosion or electrical shorts.
	\item \textbf{Mechanical Stress:} Do not apply pressure on the board or connectors when inserting/removing cables.
	\item \textbf{Transportation:} Transport in an anti-static bag inside a padded enclosure to protect against impact and ESD.
	\item \textbf{Storage:} Store the unit in a clean, dry area at room temperature when not in use.
\end{itemize}

\vspace{0.4cm}


% -------------------------------------------------------
\section{Emergency Procedures}

\begin{itemize}
	\item \textbf{Overheating:} Immediately disconnect power if the system becomes unusually hot or emits odor; inspect fan and thermal paste.
	\item \textbf{Electrical Short:} Unplug immediately, then verify continuity using a multimeter before reconnecting.
	\item \textbf{Camera or Display Failure:} Shut down system; reinsert cables after cooling; avoid hot-swapping.
	\item \textbf{Liquid Spill:} Disconnect all power; let the device dry completely (minimum 24 hours) before re-use.
	\item \textbf{Fire Risk:} Use a \textbf{CO2} or \textbf{dry chemical} extinguisher only—never water.
\end{itemize}

\vspace{0.4cm}

% -------------------------------------------------------
\section{User Responsibility}

\begin{itemize}
	\item Operators must read this safety guide before system deployment.
	\item Any modification or custom wiring outside specified limits is done at the user’s own risk.
	\item Always comply with local laboratory and institutional safety regulations.
	\item Users must ensure all firmware, drivers, and model updates are performed only through verified sources to avoid software corruption or malware.
	\item Regular documentation of system use, maintenance, and incident reports is required to maintain operational transparency and traceability.
	\item Ensure that all cables, ports, and components are handled gently during installation or inspection to prevent connector fatigue and mechanical wear.  
	\item Report any abnormal behavior such as excessive heat, noise, or repeated reboot cycles to the supervising engineer before resuming operation.  
\end{itemize}

\vspace{0.6cm}
\noindent
\textbf{Summary:}  
Adhering to these safety guidelines prevents hardware damage and personal hazards.  
Proper handling and operational awareness ensure a reliable and long-lasting Jetson Nano detection system.

% -------------------------------------------------------
% References
% -------------------------------------------------------
